% from ega4-11-fr.tex
\oldpage[IV-3]{117}
par la suite, impliquant que l'ensemble des points de platitude d'un morphisme
localement de présentation finie est ouvert.

Dans les n$^\mathrm{os}$~4 à 8, nous étudions la question de la «~descente~» de
la platitude, consistant à trouver des conditions utiles sur un morphisme de
changement de base $Y'\to Y$ (non plat en général) pour pouvoir conclure que si
$X\times_YY'$ est plat sur $Y'$, $X$ est plat sur $Y$. Ces résultats, plus
techniques que ceux des n$^\mathrm{os}$~1 à 3, sont d'une utilisation moins
fréquente dans la suite~; ils joueront cependant un rôle important dans les
techniques de construction non projectives, au chapitre suivant. Le seul
résultat des n$^\mathrm{os}$~4 à 8 utilisé dans la suite du chap.~IV est le
critère valuatif de platitude (11.8), qu'on appliquera dans (15.2).

Enfin, les n$^\mathrm{os}$~9 et 10 sont consacrés à l'étude d'une notion qui
précise, en théorie des schémas, celle de densité au sens topologique, savoir la
notion de famille de sous-préschémas \emph{schématiquement dense} dans un
préschéma donné, et notamment l'étude du comportement de cette notion par
changement de base (plat ou quelconque). Cette notion est surtout utilisée,
pour l'instant, dans l'étude des schémas en groupes.

\subsection{Ensembles de platitude (cas noethérien)}
\label{subsection:IV.11.1-fr}

\begin{theorem}[11.1.1]
\label{IV.11.1.1-fr}
Soient $Y$ un préschéma localement noethérien, $f:X\to Y$ un morphisme
localement de type fini, $\mathcal F$ un $\mathcal O_X$-Module cohérent. Alors
l'ensemble $U$ des $x\in X$ tels que $\mathcal F$ soit $f$-plat au point $x$
est ouvert dans $X$.
\end{theorem}

La question étant évidemment locale sur $X$ et $Y$, on peut supposer
$Y=\operatorname{Spec}(A)$, $X=\operatorname{Spec}(B)$, $A$ étant noethérien et
$B$ une $A$-algèbre de type fini. On a alors $\mathcal F=\widetilde M$, où $M$
est un $B$-module de type fini. Appliquons le critère
($\mathbf 0_{\mathrm{III}}$, 9.2.6)~: il suffit donc de démontrer l'assertion
suivante~:

\begin{env}[11.1.1.1]
\label{IV.11.1.1.1-fr}
\emph{Soient $A$ un anneau noethérien, $B$ une $A$-algèbre de type fini, $M$ un
$B$-module de type fini, $\mathfrak q$ un idéal premier de $B$, $\mathfrak p$
son image réciproque dans $A$. On suppose que $M_{\mathfrak q}$ soit un
$A_{\mathfrak p}$-module plat. Alors il existe $g\in B-\mathfrak q$ tel que,
pour tout idéal premier $\mathfrak q'\supset\mathfrak q$ de $B$ tel que
$g\notin\mathfrak q'$, $M_{\mathfrak q'}$ soit un
$A_{\mathfrak p'}$-module plat, en désignant par $\mathfrak p'$ l'image
réciproque de $\mathfrak q'$ dans $A$ (il revient d'ailleurs au même
($\mathbf 0_{\mathrm I}$, 6.3.1) de dire que $M_{\mathfrak q'}$ est un
$A$-module plat).}
\end{env}

Considérons pour cela $B_{\mathfrak q'}$ comme une $A$-algèbre~; on a
évidemment
$\mathfrak pB_{\mathfrak q'}\subset\mathfrak q'B_{\mathfrak q'}$. On sait
alors ($\mathbf 0_{\mathrm{III}}$, 10.2.2) que pour que $M_{\mathfrak q'}$ soit
un $A$-module plat, il faut et il suffit que
$M_{\mathfrak q'}/\mathfrak pM_{\mathfrak q'}$ soit un
$(A/\mathfrak p)$-module plat et que l'on ait
$\operatorname{Tor}_1^A(M_{\mathfrak q'},A/\mathfrak p)=0$. Or, on a
$M_{\mathfrak q'}=M\otimes_BB_{\mathfrak q'}$~; $B_{\mathfrak q'}$ étant plat
sur $B$, on a
\[
  M_{\mathfrak q'}/\mathfrak pM_{\mathfrak q'}
  =(M/\mathfrak pM)_{\mathfrak q'}
\]
et
\[
  \operatorname{Tor}_1^A(M_{\mathfrak q'},A/\mathfrak p)
  =(\operatorname{Tor}_1^A(M,A/\mathfrak p))_{\mathfrak q'}
\]
(en définissant le Tor à l'aide d'une résolution projective de
$A/\mathfrak p$)~; pour la même raison, comme on doit avoir
$g\notin\mathfrak q'$, $B_{\mathfrak q'}$ est plat sur $B_g$, donc
\[
  (M/\mathfrak pM)_{\mathfrak q'}
  =((M/\mathfrak pM)_g)_{\mathfrak q'B_g}
\]
et
\[
  (\operatorname{Tor}_1^A(M,A/\mathfrak p))_{\mathfrak q'}
  =((\operatorname{Tor}_1^A(M,A/\mathfrak p))_g)_{\mathfrak q'B_g},
\]
où dans ces formules, $M/\mathfrak pM$ et
$\operatorname{Tor}_1^A(M,A/\mathfrak p)$ sont considérés comme des
$B$-modules. Compte tenu de ($\mathbf 0_{\mathrm I}$, 6.3.2), on voit qu'on
est ramené à prouver le

\begin{lemma}[11.1.1.2]
\label{IV.11.1.1.2-fr}
Sous les conditions de (11.1.1.1), il existe $g\in B-\mathfrak q$ tel que~:
\begin{enumerate}[label=(\roman*)]
  \item $(M/\mathfrak pM)_g$ soit un $(A/\mathfrak p)$-module plat~;
  \item $(\operatorname{Tor}_1^A(M,A/\mathfrak p))_g=0$.
\end{enumerate}
\end{lemma}

